\documentclass[12pt]{article}
\usepackage{sbc-template}
\usepackage[brazil]{babel}
\usepackage[utf8]{inputenc}
\usepackage[T1]{fontenc}
\usepackage{graphicx}
\usepackage{booktabs}
\usepackage{url}
\usepackage{amsmath}

\title{Implementação e Comparação de Planejadores Clássicos para Escalonamento de Tarefas}

\author{Robson Medeiros Santos}

\address{Instituto de Ensino Superior -- ICEV\\
\email{robsonpaulista@hotmail.com}}

\begin{document}

\maketitle

\begin{resumo}
Este trabalho apresenta a modelagem de um domínio PDDL para escalonamento de tarefas em um laboratório de prototipagem, bem como a implementação e comparação de três planejadores clássicos em Python: um planejador STRIPS baseado em busca em largura/A* com heurísticas informadas, uma adaptação do Graphplan e um SATPlan com codificação CNF incremental e opção de solver externo (PySAT/Minisat22). Complementarmente, desenvolvemos um painel Streamlit para visualização gráfica dos planos e um validador automático. Os planejadores foram avaliados em três instâncias de problema, medindo-se tempo de execução, tamanho do plano e métricas específicas de cada abordagem. Os resultados mostram que todas as abordagens produzem planos válidos, mas apresentam perfis de desempenho distintos, especialmente no cenário mais complexo.
\end{resumo}

\noindent\textbf{Palavras-chave:} planejamento clássico, escalonamento de tarefas, PDDL, STRIPS, Graphplan, SATPlan

\begin{abstract}
This paper presents the modeling of a PDDL domain for task scheduling in a prototyping laboratory, as well as the implementation and comparison of three classical planners in Python: a STRIPS planner based on breadth-first search, an adaptation of Graphplan, and a SATPlan with incremental CNF encoding. The planners were evaluated on three problem instances, measuring execution time, plan length, and approach-specific metrics. Results show that all approaches yield valid plans but with distinct performance profiles, particularly in the most complex scenario.
\end{abstract}

\noindent\textbf{Keywords:} classical planning, task scheduling, PDDL, STRIPS, Graphplan, SATPlan

\section{Introdução}
Planejamento clássico fornece ferramentas formais para a geração automática de planos em sistemas determinísticos e plenamente observáveis. Nesta atividade prática, modelamos um domínio inspirado em um laboratório de fabricação digital, onde tarefas exigem operadores e máquinas específicas e respeitam restrições de precedência. O objetivo é implementar três paradigmas distintos -- STRIPS, Graphplan e SATPlan -- e comparar seu desempenho em instâncias de complexidade crescente, incorporando ainda heurísticas informadas, suporte a solver SAT externo, validação automática de planos e visualização interativa.

\section{Modelagem do Domínio}
\subsection{Elementos do domínio}
O arquivo \texttt{domain.pddl} define tipos para operadores, máquinas, recursos e tarefas, bem como predicados para disponibilidade, dependências e alocação. Utilizamos ações \texttt{start-task} e \texttt{finish-task} para capturar explicitamente a alocação/ liberação de recursos e a conclusão das tarefas.

\subsection{Restrições modeladas}
\begin{itemize}
  \item Dependências são representadas pelo predicado \texttt{(depends ?t ?p)}, garantindo que uma tarefa só possa iniciar quando todas as predecessoras estiverem concluídas.
  \item Requisitos de máquina são tratados por \texttt{(requires ?t ?m)}; quantificadores universais com igualdade asseguram que a máquina selecionada em \texttt{start-task} corresponda às exigências do problema.
  \item Uso exclusivo de recursos é modelado via predicados \texttt{available} e literais de atribuição, impedindo que um operador ou máquina execute múltiplas tarefas simultaneamente.
\end{itemize}

\section{Metodologia e Implementação}
\subsection{Planejador STRIPS}
Implementado em \texttt{planner\_strips.py}, suporta busca em largura e A* com heurísticas informadas. As heurísticas \texttt{h\_add} e \texttt{h\_max} são calculadas a partir das dependências entre tarefas, estimando o custo restante para concluir o conjunto de objetivos. O estado armazena tarefas concluídas, recursos disponíveis e tarefas em progresso. Métricas coletadas: número de ações no plano, nós expandidos e gerados, tempo total e, quando aplicável, heurística utilizada.

\subsection{Planejador Graphplan}
O arquivo \texttt{planner\_graphplan.py} constrói níveis alternados de ações e literais (positivos e negativos), computando relações de \emph{mutex} e realizando extração regressiva com memorização. As métricas incluem horizonte (número de níveis), nós expandidos durante a regressão e tempo. Literais negativos são contemplados para suportar precondições de indisponibilidade.

\subsection{Planejador SATPlan}
Em \texttt{planner\_sat.py}, codificamos o problema em CNF com horizonte incremental, facultando o uso de um solver externo (PySAT/Minisat22) ou do DPLL interno. Isso permite comparar o desempenho entre uma implementação pura e uma versão acelerada por SAT moderno. Para cada horizonte viável registramos número de variáveis, cláusulas, tamanho do plano e tempo total.

\subsection{Validação e Visualização}
O módulo \texttt{plan\_validator.py} verifica planos arbitrários simulando as ações do domínio e garante que pré-condições, alocação de recursos e metas finais sejam satisfeitas. Para análise qualitativa, \texttt{visualize\_plan.py} e o painel Streamlit geram timelines (gráficos de Gantt) com base nas ações \texttt{start-task} e \texttt{finish-task}, destacando operadores e máquinas alocadas em cada passo; as visualizações podem ser exportadas em HTML ou como imagem estática (\texttt{plotly[kaleido]}) para inclusão no relatório.

\subsection{Infraestrutura Experimental}
O script \texttt{run\_experiments.py} padroniza a execução nas três instâncias (\texttt{schedule1--3.pddl}) e consolida métricas. Todos os experimentos foram conduzidos em Python 3.11. Para os recursos opcionais (Streamlit, Plotly e PySAT) adotou-se \texttt{pip install streamlit pandas plotly python-sat[pblib,aiger]}.

\section{Resultados e Análise Comparativa}
A Tabela~\ref{tab:resultados} resume as métricas coletadas. Todos os planejadores produziram planos de mesmo tamanho, refletindo a correta modelagem do domínio. Para o STRIPS apresentamos o caso base (busca em largura); adicionalmente, as heurísticas \texttt{h\_add} e \texttt{h\_max} reduziram o volume de nós explorados na instância \texttt{schedule2} (respectivamente 24 e 65 nós expandidos, contra 72 na busca cega).

\begin{table}[ht]
\centering
\caption{Resumo dos experimentos por planejador e instância.}
\label{tab:resultados}
\small
\begin{tabular}{llrrrrr}
\toprule
Instância & Planej. & Ações & Tempo (s) & Exp. & Horiz. & Var./Cláus. \\
\midrule
schedule1 & STRIPS & 6 & 0{,}006 & 6 & -- & -- \\
schedule1 & Graph. & 6 & 0{,}020 & 7 & 6 & -- \\
schedule1 & SATPlan & 6 & 0{,}005 & -- & 6 & 134/593 \\
schedule2 & STRIPS & 8 & 0{,}007 & 72 & -- & -- \\
schedule2 & Graph. & 8 & 0{,}160 & 5 & 4 & -- \\
schedule2 & SATPlan & 8 & 250{,}703 & -- & 8 & 508/6\,784 \\
schedule3 & STRIPS & 12 & 0{,}002 & 140 & -- & -- \\
schedule3 & Graph. & 12 & 0{,}107 & 8 & 6 & -- \\
schedule3 & SATPlan & 12 & 8\,064{,}621 & -- & 12 & 743/6\,809 \\
\bottomrule
\end{tabular}
\end{table}

\subsection{Discussão}
STRIPS apresentou tempos reduzidos, e as heurísticas informadas mostraram potencial para diminuir a expansão de estados sem afetar o tamanho do plano. Graphplan manteve horizonte pequeno, obtendo planos dos mesmos tamanhos com custo moderado de extração. Já o SATPlan escalou pior quando limitado ao DPLL interno; integrações com solvers externos (PySAT/Minisat22) estão disponíveis no código e permitem reduzir drasticamente o tempo em instâncias maiores quando a biblioteca está instalada. A visualização temporal (timeline) e o validador automático auxiliaram na interpretação qualitativa e na garantia de correção dos planos produzidos por qualquer planejador.

\section{Conclusões e Trabalhos Futuros}
Os três paradigmas confirmam a corretude do domínio e produzem planos equivalentes. STRIPS é simples e eficiente para problemas de porte moderado, Graphplan equilibra horizonte e custo computacional, e SATPlan, embora geral, requer otimizações para escalabilidade. As extensões implementadas (heurísticas, PySAT, validador e visualização) elevam a reprodutibilidade e a análise qualitativa. Como trabalhos futuros sugerimos explorar heurísticas mais sofisticadas, comparar diferentes solvers SAT, enriquecer a visualização com indicadores de custo e integrar o validador a um pipeline automatizado de testes.

\section*{Agradecimentos}
Agradecemos à instituição e à coordenação da disciplina pelo suporte proporcionado.

\begin{thebibliography}{9}
\bibitem{fikes1971strips}
Richard E. Fikes e Nils J. Nilsson. STRIPS: A new approach to the application of theorem proving to problem solving. \emph{Artificial Intelligence}, 2(3–4), 1971.

\bibitem{blum1997fast}
Avrim Blum e Merrick Furst. Fast planning through planning graph analysis. \emph{Artificial Intelligence}, 90(1–2), 1997.

\bibitem{kautz1992satplan}
Henry Kautz e Bart Selman. Planning as satisfiability. \emph{ECAI}, 1992.
\end{thebibliography}

\end{document}

